\section{Definiciones y terminología}

Para el desarrollo, debemos aclarar algunas terminologías que usaremos constantemente, y son de utilidad para el lector. Dividiremos esta sección en subsecciones, para agrupar las definiciones por tema. Es importante aclarar que las definiciones están limitadas al contexto, por lo que no todas las definiciones son aplicables a todos los casos.

\subsubsection{Geodésicas}

\begin{definicion}{Geodésica}{def:geodesica}
    Una geodésica es una curva que extremiza el intervalo entre dos puntos en un espacio-tiempo. \cite{ridpath2012dictionary}
\end{definicion}

\begin{definicion}{Geodésica nula}{def:null_geodesic}
    Una geodésica nula es una curva $x^\mu(\tau)$ en una métrica $g_{\mu\nu}$ que satisface la siguiente ecuación:
    \begin{equation}
    g_{\mu\nu} \frac{dx^\mu}{d\tau} \frac{dx^\nu}{d\tau} = 0 \label{eq:null_geodesic}
    \end{equation} \cite[p. 309]{soffel2019applied}
\end{definicion}

\begin{observacion}{Rayos de luz como geodésicas nulas}{obs:rayos_de_luz_como_geodesicas_nulas}
    En relatividad general, los rayos de luz son geodésicas nulas. \cite[p. 309]{soffel2019applied}
\end{observacion}

\begin{definicion}{Geodésica de tipo tiempo}{def:timelike_geodesic}
    Una geodésica de tipo tiempo es una curva $x^\mu(\tau)$ en una métrica $g_{\mu\nu}$ que satisface la siguiente ecuación:
    \begin{equation}
    g_{\mu\nu} \frac{dx^\mu}{d\tau} \frac{dx^\nu}{d\tau} = c^2 \label{eq:timelike_geodesic}
    \end{equation} \cite[p. 169]{soffel2019applied}
\end{definicion}

\begin{observacion}{Partículas masivas como geodésicas de tipo tiempo}{obs:particulas_masivas_como_geodesicas_tipo_tiempo}
    En relatividad general, las partículas masivas siguen geodésicas de tipo tiempo. \cite[p. 169]{soffel2019applied}
\end{observacion}
